\begin{itemize}[leftmargin=*]
  \item \textbf{Mixing up mutually exclusive and independent.} Mutually exclusive events cannot happen together. Independent events do not affect each other's probabilities. These are very different ideas.
  \item \textbf{Forgetting to subtract the overlap in ``or'' questions.} If two events can happen together, the addition rule needs the intersection subtracted once.
  \item \textbf{Using constant probabilities without replacement.} Once an object is removed, the denominator usually changes.
  \item \textbf{Calling a question binomial too early.} Check all four ingredients: fixed number of trials, two outcomes per trial, constant success probability, and independence.
  \item \textbf{Confusing expectation with a guaranteed outcome.} The expected value is a long-run average, not a promise for a single trial.
  \item \textbf{Using the wrong success probability in a binomial question.} The value of $p$ must match the event you are counting as a success.
  \item \textbf{Dropping the combination factor in binomial probability.} The term $\binom{n}{r}$ counts how many different positions the successes can occupy.
  \item \textbf{Using ordered counting when order does not matter.} In committee-style questions, combinations usually keep the probability model clean.
  \item \textbf{Forgetting that probability tables must sum to $1$.} When one probability is unknown, always use this fact first.
  \item \textbf{Using a reduced sample space incorrectly in conditional probability.} Once the condition is given, the denominator should reflect that condition only.
\end{itemize}

\noindent A strong exam habit is to write one short line naming the method before calculating. For example:
\begin{itemize}[leftmargin=*]
  \item ``Use complement since `at least one' is awkward directly.''
  \item ``Use conditional probability because the sample space has changed.''
  \item ``This is binomial with $n=5$ and $p=0.3$.''
\end{itemize}

\noindent That brief method statement often prevents a correct formula from being used in the wrong setting.
